\documentclass[10pt,a4paper]{amsart}
\usepackage[margin=19mm]{geometry}
\usepackage{amsmath,amssymb,mathtools}
\usepackage[scr=rsfs,cal=euler]{mathalfa}
\usepackage{xfrac}
\usepackage[unicode,hidelinks,bookmarks=false]{hyperref}
\newcommand{\ie}{i.e.\ }
\newcommand{\cf}{cf.\ }
\newcommand{\resp}{respectively\ }
\newcommand{\Subsection}{\subsection}
\providecommand{\nobreakdash}{\nobreak\hbox{-}}
\setlength{\emergencystretch}{2em}
\multlinegap0pt
\usepackage{amssymb}
\newtheorem{lemma}{Lemma}
\newtheorem{proposition}{Proposition}
\newcommand{\PP}{{\mathbb P}}
\newcommand{\calS}{{\mathcal S}}
\title{Rational points on smooth surfaces in $\PP^3$ over finite fields}
\author{Yves Aubry, Jos\'e Felipe Voloch}
\date{Source: arXiv:2605.09166v1 (CC BY 4.0)}
\begin{document}
\maketitle

\noindent\emph{Source and licence.} Rational points on smooth surfaces in $\PP^3$ over finite fields. Version arXiv:2605.09166v1 (CC BY 4.0), distributed by arXiv under the Creative Commons Attribution 4.0 International licence, http://creativecommons.org/licenses/by/4.0/, as recorded in the arXiv licence metadata for this exact version and shown on the abstract page https://arxiv.org/abs/2605.09166v1. Full licence notice: \texttt{licenses/arxiv-2605.09166-CC-BY-4.0.txt}.\par\smallskip

\noindent\textit{Editorial scope.} Counted unit: the article's proposition proposition (source0.tex lines 424--430). The article's complete original proof of it is delivered with it (source0.tex lines 432--517, one \texttt{proof} environment of 86 lines). No mathematics is written, repaired or re-derived: the statement and the proof are the article's own lines, in the article's own notation, and no retained line is substituted or re-flowed.  Retained uncounted as the source-local context the proof needs: lines 218--221; lines 364--367.  Nothing was omitted from any retained span, so the package carries no omission record.  No retained line contains a citation, so the wrapper carries no bibliography and no bibliography entry.  No retained line contains a figure environment, an image-inclusion command or drawing code, so no image is delivered and the package declares no asset dependency.  Editorial formatting only, in addition: an AMS \texttt{amsart} preamble that restates the article's own declarations and macros used by the retained lines, namely the environments \texttt{lemma}, \texttt{proposition} on separate counters, together with the standard packages those lines need.  The retained environments print in this document as Lemma 1, Proposition 1 in order of appearance, whereas the article prints the same items with the numbering of its own full document.  The complete native source of the article as distributed in the arXiv e-print for this exact version is bundled unchanged as \texttt{sources/200202/source0.tex}.
\begin{lemma}\label{Fondamental_Lemma}
    Let $\calS$ be a smooth surface in $\PP^3$ of degree $d$ and $\calS_i$ surfaces in $\PP^3$ of degree $d_i, i=1,2$. Assume that $\calS,\calS_1,\calS_2$ intersect in a zero-dimensional scheme $\Gamma$ and lines $L_1,\ldots,L_m$ with multiplicity one. Then
    $$2m \le \deg \Gamma - (dd_1d_2 -m(d+d_1+d_2)) \le m(m+1).$$
\end{lemma}

\begin{proposition}\label{Number_Lines}
If $-1$ is a $d$-th power modulo $p$ then the  surface $S$ contains $3d^2$ lines defined over ${\mathbb F}_p$ and
all the lines in the intersection $X:=S\cap S_1 \cap S_2$ have multiplicity one.
\end{proposition}

\begin{proposition}
The number of rational points of the surface $S$ in ${\mathbb P}^3$ defined over ${\mathbb F}_p$  by 
$x^{(p-1)/2} + y^{(p-1)/2} -z^{(p-1)/2} -w^{(p-1)/2}=0$ is given by:

$$\sharp S({\mathbb F}_p)=\frac{3}{8}(p^3-3p^2+11p-9).$$
This is an example of a surface such that $S\cap S_1\cap S_2$ does not contain isolated rational points.
\end{proposition}

\begin{proof}
Let us remark that if $a\in{\mathbb F}_p$ then, $a^{(p-1)/2}=0$ if and only if $a=0$ and  it is equal to 1 if $a$ is a quadratic residue modulo $p$ and is equal to $-1$ otherwise. And it is well known that the number of nonzero quadratic residues, as the number of  quadratic nonresidues modulo $p$ is equal to $(p-1)/2$. We begin by counting the number of affine rational points on $S$. When none of the coordinates are zero, we have 6 choices for the positions of the $+1$, so $6(\frac{p-1}{2})^4$ solutions. If two coordinates are zero, then we have also 6 choices for the zero coordinates and then two possibilities for the choices of $+1$ or $-1$, thus $6\times 2 \times (\frac{p-1}{2})^2$ solutions. Adding the point with all coordinate zero, we obtain the number $N_{\rm aff}$ of affine rational points on $S$ which is equal to 
$6(\frac{p-1}{2})^4+12 (\frac{p-1}{2})^2+1$ and then $\sharp S({\mathbb F}_p) =\frac{N_{\rm aff}-1}{p-1}$ which gives the result.

Moreover, let us calculate the degree of $\Gamma$, the zero-dimensional scheme appearing in the intersection 
$X:=S\cap S_1\cap S_2$.
By Proposition \ref{Number_Lines} we have that the lines  
in $X$ have multiplicity one, and so we can apply
Lemma \ref{Fondamental_Lemma} and we get
$\deg\Gamma=dd_1d_2 - m(d+d_1+d_2)+2m+ \sum_{i \ne j} L_i\cdot L_j$
with, here, $d=(p-1)/2, d_1=3(p-1)/2$ and $d_2=5(p-1)/2$. Let us scrutinize now the intersections of the lines $L_i$.
Since $-1$ is clearly a $(p-1)/2$-th power modulo $p$, one can use
 Proposition \ref{Number_Lines} to asserts that the lines $L_i$ contained in  $X$
 are of three types:
\begin{equation*}
\left\{
\begin{matrix}
x = ay   \\
z = bw \cr
\end{matrix}
\right.
\qquad
\left\{
\begin{matrix}
x = a'z   \\
y = b'w\cr
\end{matrix}
\right. 
\qquad
\left\{
\begin{matrix}
x = a' w   \\
y = b'z \cr
\end{matrix}
\right. 
\end{equation*}
where $a,b$ are quadratic non-residues modulo $p$ and $a',b'$ are (nonzero) quadratic residues modulo $p$.
Thus we obtain, according to Proposition \ref{Number_Lines}, that  the number of lines is $m=3(\frac{p-1}{2})^2$.

If $L_1$ and $L_2$ are two lines of same type given, for example, by equations
\begin{equation*}
\left\{
\begin{matrix}
x = a_1y   \\
z = b_1w \cr
\end{matrix}
\right.
\qquad
\left\{
\begin{matrix}
x = a_2y   \\
z = b_2w\cr
\end{matrix}
\right. 
\end{equation*}
then $L_1.L_2=1$ if and only if $a_1=a_2$ or $b_1=b_2$.
Hence the contribution will be in the case of lines $L_i$ and $L_j$ of same type:
$$\sum_{i\not=j, {\rm same\ type}}L_i.L_j=3\left(\frac{p-1}{2}\right)^2\left((\frac{p-1}{2}-1)\times 2\right)=3\left(\frac{p-1}{2}\right)^2(p-3).$$

Moreover, if $L_1$ and $L_2$ are two lines of different type given, for example, by equations
\begin{equation*}
\left\{
\begin{matrix}
x = ay   \\
z = bw \cr
\end{matrix}
\right.
\qquad
\left\{
\begin{matrix}
x = a'z   \\
y = b'w\cr
\end{matrix}
\right. 
\end{equation*}
then $L_1.L_2=1$ if and only if $ab'=a'b$.
Hence the contribution will be in the case of lines $L_i$ and $L_j$ of different type:
$$\sum_{i\not=j, {\rm different\  type}}L_i.L_j=6\left(\frac{p-1}{2}\right)^3.$$
Finally, we obtain:
$$\sum_{i\not=j}L_i.L_j=6\left(\frac{p-1}{2}\right)^2(p-2).$$

One can compute now 
$\deg\Gamma= dd_1d_2 - m(d+d_1+d_2)+2m+ \sum_{i \ne j} L_i\cdot L_j$ which gives
$$\deg\Gamma= 15\left(\frac{p-1}{2}\right)^3-3\left(\frac{p-1}{2}\right)^2\left(9\left(\frac{p-1}{2}\right)\right)+6\left(\frac{p-1}{2}\right)^2+6\left(\frac{p-1}{2}\right)^2(p-2)=0.$$
Thus $S\cap S_1\cap S_2$ does not contain isolated rational points.
\end{proof}

\end{document}
